\section{Kontaktlogg}
\label{app:kontaktlogg}

Dette vedlegget inneholder en systematisk oversikt over kvalitative intervjuer og samtaler gjennomført i forbindelse med mulighetsanalysen. Formålet med loggen er å dokumentere innsamlet informasjon på en måte som gjør den lett å tolke og bruke videre i utviklingen av forretningsmodellen.

Kontaktpersonene er organisert i fire kategorier basert på rolle: potensielle kunder, brukere, samarbeidspartnere og eksperter. Denne inndelingen gjør det enklere å analysere innsikter ut fra forskjellige perspektiver og identifisere hvilke behov og utfordringer som er mest relevante for problemstillingen.

Som påpekt av \textcite[57]{torvatn_teknologiledelse_2024}, er aktiv kontakt med relevante aktører – herunder kunder, leverandører, produsenter og distributører – avgjørende for å hente informasjon som støtter utforming av forretningsmodellen. Kontaktloggen følger derfor en struktur som registrerer hver interaksjon, inkludert dato, ansvarlig student og kontaktinformasjon, slik at innsamlet data kan brukes systematisk i analysen.

Ved å dokumentere alle telefonsamtaler, e-poster og møter sikrer loggen at viktig informasjon ikke går tapt. Samtidig gjør den det mulig å tolke og bruke materialet systematisk, slik at både studentgruppen og andre kan dra nytte av innsamlet data i analysen og trekke ut sentrale funn på tvers av kontaktene \parencite[57]{torvatn_teknologiledelse_2024}.

Intervjuene ble gjennomført som semistrukturerte dybdeintervjuer. For hver målgruppe ble det utarbeidet en intervjuguide med faste tema og åpne spørsmål for å sikre både sammenlignbarhet og fleksibilitet. Dette gjorde det mulig å utforske uforutsette innsikter samtidig som sentrale problemstillinger ble dekket i alle samtaler. Samtalene ble dokumentert gjennom notater og oppsummert umiddelbart etter gjennomføring.

\subsection{Kunder}

\subsubsection{Kontakt 1 – Ola Nordmann}

\textbf{Bransje/stilling:} Bla bla bla \\
\textbf{Kontaktinfo:} E-post/tlf/etc. \\
\textbf{Dato:} 12.01.2026 \\
\textbf{Ansvarlig fra studentgruppen:} Én av oss her

\paragraph{Sentrale funn}
Bla bla bla bla...

\subsubsection{Kontakt 2 – Kari Hansen}

\textbf{Bransje/stilling:} \\
\textbf{Kontaktinfo:} \\
\textbf{Dato:} \\
\textbf{Ansvarlig fra studentgruppen:}

\paragraph{Sentrale funn}

\subsection{Brukere}

\subsubsection{Kontakt 3 – Andreas Berg}

\textbf{Bransje/stilling:} \\
\textbf{Kontaktinfo:} \\
\textbf{Dato:}  \\
\textbf{Ansvarlig fra studentgruppen:} 

\paragraph{Sentrale funn}

\subsection{Samarbeidspartnere}

\subsubsection{Kontakt 4 – Ingrid Lunde}

\textbf{Bransje/stilling:}  \\
\textbf{Kontaktinfo:}  \\
\textbf{Dato:} \\
\textbf{Ansvarlig fra studentgruppen:}

\paragraph{Sentrale funn}

\subsection{Eksperter}

\subsubsection{Kontakt 5 – Albert Lau}
\label{kontakt:lau}

\textbf{Fagområde/stilling:} Førsteamenuensis NTNU \\
\textbf{Kontaktinfo:} albert.lau@ntnu.no, +47 942 58 319 \\
\textbf{Dato:} 04.03.2026 \\
\textbf{Ansvarlig fra studentgruppen:} Sigurd Riseth

\paragraph{Sentrale funn} 

Albert Lau bekrefter at jernbaneinfrastruktur er generelt mer sårbar enn veiinfrastruktur, særlig i Norge med mye enkeltspor --- en feil kan stoppe trafikk på hele strekninger. Han trekker frem bruer og sporveksler som de mest kritiske objektene å overvåke, der spesielt sporveksler i trafikkintensive knutepunkter (som Oslo S og Alnabru) har stor operasjonell betydning. Viktigste årsaker til tilstandsforringelse er alder og utmattingsskader (fatigue).

Dagens inspeksjonsmetoder er manuelle inspeksjoner og målevognen Roger 1000. Den største svakheten ligger ikke i datainnsamlingen, men i dataanalysen – data sammenlignes fortsatt i stor grad manuelt mot grenseverdier. Lau peker på AI-drevet analyse som det største forbedringspotensiale, og hans eget forskningsmiljø jobber med nettopp dette.

På spørsmål om barrierer for autonom droneinspeksjon mener han de organisatoriske og regulatoriske utfordringene er større enn de tekniske. Organisatorisk handler det om motstand og skepsis mot ny teknologi. Regulatorisk er datadeling og AI-regulering i Europa særlige utfordringer.

Betalingsvilje for automatiserte overvåkingsløsninger er til stede, men forutsetter dokumentert effekt – gjerne med referanser fra andre europeiske land. Lau nevner at Bane NOR allerede kjøper teknologi fra bl.a. Østerrike, som et eksempel på at slik vilje finnes.

\subsubsection{Kontakt 6 – Emad Samuel Malki Ebeid}

\textbf{Fagområde/stilling:} Professor på Syddansk Universitet, Prosjekt koordinator for Drones4Safety.eu \\
\textbf{Kontaktinfo:} esme@mmmi.sdu.dk \\
\textbf{Dato:} 12.02.2026 \\
\textbf{Ansvarlig fra studentgruppen:} Sigurd Riseth

\paragraph{Sentrale funn} 

Emad bekrefter at Drones4Safety utviklet et system av autonome, selvladende droner for inspeksjon av broer og jernbane. Det mest sentrale resultatet er verdens første vellykkede landing og lading på en høyspentlinje, demonstrert ved Aarhus Letbane i 2022 – noe Emad omtaler som potensielt revolusjonerende for dronebransjen, ettersom typisk flygetid i dag er begrenset til 10–15 minutter.

Prosjektet utviklet sværmalgoritmer for samordnet flerdroneoperasjon, samt AI-algoritmer for automatisk deteksjon av infrastrukturfeil fra dronebilder via skybasert analyse. Systemet ble testet på Asti-broen i Italia og Siemens' jernbanetestsenter i Tyskland.

Emad peker på regulatoriske forhold som den viktigste barrieren for kommersiell utrulling – særlig lovgivning rundt autonom luftfart – og anbefaler å ta tidlig kontakt med relevante myndigheter som EASA og Luftfartstilsynet. Han understreker videre at AI-modeller bør bygge på eksisterende treningsdata fra prosjekter som Drones4Safety fremfor å utvikles fra bunnen av. Han anslår at kommersielt tilgjengelige selvladende droner kan være realitet innen 3–4 år, og estimater indikerer et kostnadsbesparelsespotensial på €15 mrd. årlig for inspeksjon av europeisk transportinfrastruktur.

\subsubsection{Kontakt 7 – Inge Hoff}

\textbf{Fagområde/stilling:} Professor på NTNU \\
\textbf{Kontaktinfo:} inge.hoff@ntnu.no, +47 735 94 731, +47 934 26 463 \\
\textbf{Dato:}  \\
\textbf{Ansvarlig fra studentgruppen:} Sigurd Riseth

\paragraph{Sentrale funn} 
